\documentclass[../../main.tex]{subfiles}% 注意这里的文件路径不能用 ./main.tex ,否则用latexmk编译子文件会报错
\graphicspath{{\subfix{./image/}}} % 指定图片目录,后续可以直接使用图片文件名
% 注意这里的文件路径不能用 ../../image/ ,否则用latexmk编译子文件会报错

% 例如:
% \begin{figure}[H]
% \centering
% \includegraphics[scale=0.3]{图.png}
% \caption{}
% \label{figure:图}
% \end{figure}
% 注意:上述\label{}一定要放在\caption{}之后,否则引用图片序号会只会显示??.

\begin{document}

\section{自然标架的运动公式}
正则参数曲面的自然标架为$\{\bm{r};\bm{r}_u,\bm{r}_v,\bm{n}\}$，其运动公式为
\begin{align*}
&\bm{r}_{uu}=\Gamma^1_{11}\bm{r}_u+\Gamma^2_{11}\bm{r}_v+L\bm{n},\\
&\bm{r}_{uv}=\Gamma^1_{12}\bm{r}_u+\Gamma^2_{12}\bm{r}_v+M\bm{n},\\
&\bm{r}_{vv}=\Gamma^1_{22}\bm{r}_u+\Gamma^2_{22}\bm{r}_v+N\bm{n},\\
&\bm{n}_u=-\dfrac{LF-ME}{EG-F^2}\bm{r}_u-\dfrac{MF-NE}{EG-F^2}\bm{r}_v,\\
&\bm{n}_v=-\dfrac{LM-MF}{EG-F^2}\bm{r}_u-\dfrac{MN-NF}{EG-F^2}\bm{r}_v,
\end{align*}
其中$\Gamma^i_{jk}$是Christoffel记号，由第一类基本量唯一确定。

\section{Gauss方程和Codazzi方程}
\begin{theorem}
曲面的第一、第二基本量满足Gauss方程
\begin{align*}
&(\Gamma^1_{22})_u-(\Gamma^1_{12})_v+\Gamma^2_{22}\Gamma^1_{12}-\Gamma^1_{12}\Gamma^2_{12}+\Gamma^1_{22}\Gamma^1_{11}-\Gamma^1_{11}\Gamma^1_{22}=\dfrac{EM-LF}{EG-F^2},\\
&(\Gamma^2_{11})_v-(\Gamma^2_{12})_u+\Gamma^1_{11}\Gamma^2_{12}-\Gamma^2_{12}\Gamma^1_{12}+\Gamma^2_{11}\Gamma^2_{22}-\Gamma^2_{12}\Gamma^2_{11}=\dfrac{FN-MG}{EG-F^2},
\end{align*}
和Codazzi方程
\begin{align*}
L_v-M_u&=L\Gamma^1_{12}+M(\Gamma^2_{12}-\Gamma^1_{11})-N\Gamma^2_{11},\\
M_v-N_u&=L\Gamma^1_{22}+M(\Gamma^2_{22}-\Gamma^1_{12})-N\Gamma^2_{12}.
\end{align*}
\end{theorem}

\begin{corollary}
高斯曲率$K$仅由第一基本形式决定，是曲面的内蕴不变量。
\end{corollary}

\section{曲面论基本定理}
\begin{theorem}
设$D\subset\mathbb{R}^2$是单连通区域，$E,F,G$和$L,M,N$是$D$上的光滑函数，满足
\begin{enumerate}[(1)]
\item $E>0,\ G>0,\ EG-F^2>0$；
\item 满足Gauss方程和Codazzi方程，
\end{enumerate}
则在$\mathbb{E}^3$中存在唯一的正则参数曲面（不计刚体运动），以$E,F,G$为第一类基本量，$L,M,N$为第二类基本量。
\end{theorem}

\begin{proof}
由自然标架的运动公式构成一阶线性偏微分方程组，由可积条件（Gauss-Codazzi方程）可知方程组可积，解存在且唯一，差一个刚体运动。

\end{proof}

\section{常高斯曲率曲面}
\begin{theorem}
常正高斯曲率曲面局部等距于球面，常零高斯曲率曲面局部等距于平面，常负高斯曲率曲面局部等距于伪球面。
\end{theorem}

\begin{example}
伪球面的参数方程为
\[
\bm{r}(u,v)=\left(a\sin u\cos v,a\sin u\sin v,a\left(\cos u+\ln\tan\frac{u}{2}\right)\right),
\]
其高斯曲率$K=-\dfrac{1}{a^2}$。
\end{example}

\section{极小曲面}
\begin{definition}
平均曲率$H\equiv0$的曲面称为**极小曲面**。
\end{definition}

\begin{example}
正螺旋面、悬链面是典型的极小曲面，且正螺旋面与悬链面等距。
\end{example}

\begin{theorem}
曲面为极小曲面的充要条件是曲面上任意邻域的面积为临界值。
\end{theorem}

\section{曲面上正交参数曲线网的存在性}
\begin{lemma}
设$f(u,v),g(u,v)$是区域$D\subset\mathbb{E}^2$上不同时为零的连续可微函数，则对任意$(u_0,v_0)\in D$，存在邻域$U\subset D$及非零连续函数$\lambda(u,v)$，使得$\lambda(u,v)$是一次微分式$\omega=f(u,v)\mathrm{d}u+g(u,v)\mathrm{d}v$的积分因子，即在$U$上存在连续可微函数$k(u,v)$，使得
\[
\lambda(u,v)\left(f(u,v)\mathrm{d}u+g(u,v)\mathrm{d}v\right)=\mathrm{d}k(u,v).
\]
\end{lemma}

\begin{theorem}
假定在正则参数曲面$S:\bm{r}=\bm{r}(u,v)$上有两个处处线性无关的连续可微切向量场$\bm{a}(u,v),\bm{b}(u,v)$，则对任意$p\in S$，存在$p$的邻域$U\subset S$及新参数系$(\tilde{u},\tilde{v})$，使得新参数曲线的切向量$\bm{r}_{\tilde{u}},\bm{r}_{\tilde{v}}$分别与$\bm{a},\bm{b}$平行。
\end{theorem}

\begin{proof}
设$\bm{a}=a_1\bm{r}_u+a_2\bm{r}_v$，$\bm{b}=b_1\bm{r}_u+b_2\bm{r}_v$，其中$a_i,b_i$连续可微，且行列式
\[
A=\begin{vmatrix}a_1&a_2\\b_1&b_2\end{vmatrix}\neq0.
\]
若存在参数变换使得$\bm{r}_{\tilde{u}}\parallel\bm{a}$，$\bm{r}_{\tilde{v}}\parallel\bm{b}$，则$\mathrm{d}\tilde{u},\mathrm{d}\tilde{v}$应满足
\[
\begin{cases}
a_1\mathrm{d}u+a_2\mathrm{d}v=0,\\
b_1\mathrm{d}u+b_2\mathrm{d}v=0.
\end{cases}
\]
由引理，存在积分因子$\lambda,\mu$使得$\lambda(a_1\mathrm{d}u+a_2\mathrm{d}v)=\mathrm{d}\tilde{u}$，$\mu(b_1\mathrm{d}u+b_2\mathrm{d}v)=\mathrm{d}\tilde{v}$，对应的参数变换即为所求。

\end{proof}

\begin{corollary}
正则参数曲面的每一点邻域内必存在正交参数曲线网，使得$F\equiv0$。
\end{corollary}

\section{保长对应和保角对应}
\begin{definition}
若曲面$S_1$到$S_2$的连续可微映射$\sigma$满足$\sigma^*\mathrm{I}_2=\mathrm{I}_1$，即保持切向量长度不变，则称$\sigma$为保长对应。
\end{definition}

\begin{definition}
若曲面$S_1$到$S_2$的连续可微映射$\sigma$满足$\sigma^*\mathrm{I}_2=\lambda^2\mathrm{I}_1$（$\lambda>0$为比例因子），即保持切向量夹角不变，则称$\sigma$为保角对应。
\end{definition}

\begin{theorem}
映射$\sigma:S_1\to S_2$为保角对应的充要条件是存在正函数$\lambda$，使得
\[
\begin{cases}
\tilde{E}=\lambda^2 E,\\
\tilde{F}=\lambda^2 F,\\
\tilde{G}=\lambda^2 G.
\end{cases}
\]
\end{theorem}

\begin{theorem}
任意正则参数曲面的每一点邻域，都可与平面区域建立保角对应。
\end{theorem}

\begin{example}
球面与平面之间的球极投影是保角对应。
\end{example}

\section{可展曲面}
\begin{definition}
单参数直线族生成的曲面称为直纹面，参数方程为$\bm{r}(u,v)=\bm{a}(u)+v\bm{l}(u)$，其中$\bm{a}(u)$为准线，$\bm{l}(u)$为直母线方向向量。
\end{definition}

\begin{definition}
高斯曲率恒为零的直纹面称为可展曲面，可展曲面局部上是平面、柱面、锥面或切线面。
\end{definition}

\begin{theorem}
直纹面$\bm{r}(u,v)=\bm{a}(u)+v\bm{l}(u)$为可展曲面的充要条件是混合积$(\bm{a}'(u),\bm{l}(u),\bm{l}'(u))\equiv0$。
\end{theorem}

\begin{proof}
直纹面的第一基本量满足$EG-F^2=(\bm{a}'+v\bm{l}')\times\bm{l}|^2$，第二基本量满足$LN-M^2\equiv0$等价于混合积为零。

\end{proof}

\begin{remark}
可展曲面能够无撕裂、无伸缩地展开为平面，即与平面存在保长对应。
\end{remark}



\end{document}